%%%%%%%%%%%%%%%%%%%%%%%%%%%%%%%%%%%%%%%%%%%%%%%%%%%%%%%%
%  yinyong.tex — 引用与参考文献
%  演示 kaobook + biblatex 引用体系的核心用法
%%%%%%%%%%%%%%%%%%%%%%%%%%%%%%%%%%%%%%%%%%%%%%%%%%%%%%%%

\setchapterpreamble[u]{\margintoc}
\chapter{引用与参考文献}
\labch{yinyong}

在学术写作中，引用与参考文献是至关重要的一环。本章系统介绍本模板基于
\Package{biblatex}（通过 \Package{kaobiblioCJKsc} 加载）和
\Package{kaorefsCJKsc} 所提供的中文引用体系，涵盖文献引用、图表交叉引用、
以及 \Package{cleveref} 智能引用三大板块。

%----------------------------------------------------------------------------------------

\section{文献引用基础}

\index{引用!文献引用}
本模板使用 \Package{biblatex} 作为参考文献引擎，后端为 \texttt{biber}。
所有参考文献命令均继承自 \Package{biblatex}，支持多种引用风格。

\subsection{基本引用命令}

\index{引用!parencite@\texttt{\textbackslash parencite}}
\index{引用!textcite@\texttt{\textbackslash textcite}}
\index{引用!citeauthor@\texttt{\textbackslash citeauthor}}
\index{引用!citeyear@\texttt{\textbackslash citeyear}}
\index{引用!cite@\texttt{\textbackslash cite}}

以下是 \Package{biblatex} 最常用的引用命令及其示例：

\begin{table}[h]
\caption{常用文献引用命令一览}
\labtab{yinyong-commands}
\begin{tabular}{@{}llp{6cm}@{}}
\toprule
命令 & 输出格式 & 说明 \\
\midrule
\lstinline|\cite{key}| & [1] 或 (作者, 年份) & 基本引用，取决于样式 \\
\lstinline|\parencite{key}| & (作者, 年份) & 带括号引用 \\
\lstinline|\textcite{key}| & 作者 (年份) & 文中引用，作者名融入句子 \\
\lstinline|\citeauthor{key}| & 作者 & 仅引用作者名 \\
\lstinline|\citeyear{key}| & 年份 & 仅引用年份 \\
\lstinline|\citeitle{key}| & 标题 & 仅引用文献标题 \\
\lstinline|\footcite{key}| & 脚注引用 & 在脚注中给出完整信息 \\
\bottomrule
\end{tabular}
\end{table}

\marginnote[-2cm]{在实际编译时，你需要将上表中的 \lstinline|key| 替换为%
\lstinline|.bib| 文件中的文献条目键值。}

例如，文中引用可以写作：\lstinline|\textcite{Author2024} 提出了一种新方法|，
这将生成类似「张三 (2024) 提出了一种新方法」的输出。

\subsection{biblatex 引用样式}

\index{引用!样式}
\Package{biblatex} 支持多种引用和参考文献样式，可通过加载
\Package{kaobiblioCJKsc} 时传入选项来指定。常用样式包括：

\begin{itemize}
    \item \texttt{style=numeric} —— 数字编号，如 [1], [2], \ldots{}
    \item \texttt{style=authoryear} —— 作者-年份，如 (张三, 2024)
    \item \texttt{style=authortitle} —— 作者-标题
    \item \texttt{style=verbose} —— 详尽脚注式
    \item \texttt{style=reading} —— 阅读型，适合人文学科
    \item \texttt{style=chicago-authordate} —— Chicago 格式
\end{itemize}

加载方式示例：
\begin{lstlisting}
% 使用数字编号样式
\usepackage[style=numeric]{kaobiblioCJKsc}
% 或使用作者-年份样式
\usepackage[style=authoryear]{kaobiblioCJKsc}
\end{lstlisting}

\index{引用!sidecite@\texttt{\textbackslash sidecite}}
此外，本模板提供了 \lstinline|\sidecite{key}| 命令，它会将引用信息和
文献标题显示在页边注中——这是 kaobook 体系的标志性特征。

默认的侧边引用格式由 \lstinline|\formatmargincitation| 命令控制：
\begin{lstlisting}
\newcommand{\formatmargincitation}[1]{%
    \parencite{#1}: \citeauthor*{#1} (\citeyear{#1}), \citetitle{#1}%
}
\end{lstlisting}

你可以通过 \lstinline|\renewcommand{\formatmargincitation}| 来自定义
侧边引用的显示格式。

%----------------------------------------------------------------------------------------

\section{图表与定理交叉引用}

\index{引用!交叉引用}
\index{引用!label@\texttt{\textbackslash label}}
\index{引用!ref@\texttt{\textbackslash ref}}

\subsection{基本引用命令}

\Package{kaorefsCJKsc} 提供了一套层次化的引用命令，用于引用章节、
图表、公式和定理等元素。每种元素都有四种引用深度：

\begin{table}[h]
\caption{交叉引用命令体系}
\labtab{kaorefs-commands}
\begin{tabular}{@{}lll@{}}
\toprule
引用深度 & 命令前缀 & 示例（用于章节）\\
\midrule
纯编号 & \texttt{\textbackslash ref} & \lstinline|\refch{yinyong}| $\rightarrow$ 章~\refch{yinyong} \\
编号 + 页码 & \texttt{\textbackslash vref} & \lstinline|\vrefch{yinyong}| \\
编号 + 标题 & \texttt{\textbackslash nref} & \lstinline|\nrefch{yinyong}| \\
完整（编号+标题+页码） & \texttt{\textbackslash fref} & \lstinline|\frefch{yinyong}| \\
\bottomrule
\end{tabular}
\end{table}

\subsection{标签规约}

\index{引用!labch@\texttt{\textbackslash labch}}
为保持一致性，\Package{kaorefsCJKsc} 为每种元素定义了专用标签命令，
它们自动添加前缀以避免冲突：

\begin{multicols}{2}
\begin{itemize}
    \item \lstinline|\labch{name}| $\rightarrow$ \lstinline|ch:name| （章节）
    \item \lstinline|\labsec{name}| $\rightarrow$ \lstinline|sec:name| （节）
    \item \lstinline|\labfig{name}| $\rightarrow$ \lstinline|fig:name| （图）
    \item \lstinline|\labtab{name}| $\rightarrow$ \lstinline|tab:name| （表）
    \item \lstinline|\labeq{name}| $\rightarrow$ \lstinline|eq:name| （公式）
    \item \lstinline|\labthm{name}| $\rightarrow$ \lstinline|thm:name| （定理）
    \item \lstinline|\labdef{name}| $\rightarrow$ \lstinline|def:name| （定义）
\end{itemize}
\end{multicols}

相应的引用命令也有对应后缀，例如：
\begin{itemize}
    \item 引用图：\lstinline|\reffig{name}|, \lstinline|\vreffig{name}|
    \item 引用表：\lstinline|\reftab{name}|, \lstinline|\vreftab{name}|
    \item 引用公式：\lstinline|\refeq{name}|, \lstinline|\vrefeq{name}|
\end{itemize}

\subsection{图表示例}

\index{引用!图引用}
以下是一个带标签的示意表格：

\begin{table}[h]
\centering
\caption{引用命令速查——kaorefs 体系}
\labtab{ref-quick}
\begin{tabular}{@{}lll@{}}
\toprule
元素 & 标签命令 & 引用命令（纯编号）\\
\midrule
章 & \lstinline|\labch{}| & \lstinline|\refch{}| \\
节 & \lstinline|\labsec{}| & \lstinline|\refsec{}| \\
图 & \lstinline|\labfig{}| & \lstinline|\reffig{}| \\
表 & \lstinline|\labtab{}| & \lstinline|\reftab{}| \\
公式 & \lstinline|\labeq{}| & \lstinline|\refeq{}| \\
定理 & \lstinline|\labthm{}| & \lstinline|\refthm{}| \\
\bottomrule
\end{tabular}
\end{table}

引用上表：\reftab{ref-quick}（这里展示的是纯编号引用）。
使用 \lstinline|\vreftab{ref-quick}| 还会附加页码信息。

%----------------------------------------------------------------------------------------

\section{cleveref 智能引用}

\index{引用!cleveref@\texttt{cleveref}}

\Package{cleveref} 是 \LaTeX{} 社区广泛使用的智能引用宏包，它能自动判断
被引用对象的类型并生成对应的文字前缀。中文模板已通过
\Package{kaorefsCJKsc} 自动加载 \Package{cleveref}。

\subsection{cleveref 基本用法}

\index{引用!cref@\texttt{\textbackslash cref}}
\index{引用!Cref@\texttt{\textbackslash Cref}}

\Package{cleveref} 的核心命令是 \lstinline|\cref{}| 和 \lstinline|\Cref{}|：

\begin{itemize}
    \item \lstinline|\cref{label1,label2,...}| —— 小写开头，如「图~1 和 2」
    \item \lstinline|\Cref{label1,label2,...}| —— 大写开头（句首使用），如「Figure~1 and 2」
\end{itemize}

\marginnote{中文环境中 \texttt{\textbackslash Cref} 与 \texttt{\textbackslash cref} %
的区别主要体现在英文标题的句首大小写上；中文本身没有大小写问题。}

\Package{cleveref} 优势在于它能自动识别标签类型，并处理多标签的并列关系。
例如，\lstinline|\cref{fig:A,fig:B,fig:C}| 将输出「图~1--3」，
而 \lstinline|\cref{fig:A,tab:X}| 将输出「图~1 和 表~1」。

\subsection{cleveref 与 kaorefs 的配合}

\index{引用!kaorefs 与 cleveref@kaorefs 与 cleveref}
由于 \Package{kaorefsCJKsc} 的标签命令（如 \lstinline|\labfig{}|）本质上
就是 \lstinline|\label{}|，因此 \Package{cleveref} 的 \lstinline|\cref{}|
可以直接作用于这些标签。

\begin{lstlisting}
% 使用 kaorefs 标签 + cleveref 引用
\labfig{sample} % 等价于 \label{fig:sample}
% ... 图的内容 ...
如 \cref{fig:sample} 所示，这个方法是有效的。
\end{lstlisting}

\subsection{中文本地化}

\index{引用!中文本地化}
在中文模板中，\Package{kaorefsCJKsc} 已通过 \lstinline|config.tex| 设置了
中文本地化字符串（如「图」「表」「章」等）。如需自定义 \Package{cleveref}
的类型名称，可使用：

\begin{lstlisting}
\crefname{figure}{图}{图}
\Crefname{figure}{图}{图}
\crefname{table}{表}{表}
\Crefname{table}{表}{表}
\end{lstlisting}

%----------------------------------------------------------------------------------------

\section{编译参考文献}

\index{引用!编译}
\index{编译!biber@biber}
使用 \Package{biblatex} 时，需要额外运行 \texttt{biber} 来处理参考文献数据库。
典型的编译流程如下：

\begin{lstlisting}[style=kaolstplain]
xelatex main
biber main
xelatex main
xelatex main   % 第三次编译以确保交叉引用稳定
\end{lstlisting}

或者使用 \texttt{latexmk} 自动化编译：

\begin{lstlisting}[style=kaolstplain]
latexmk -xelatex main
\end{lstlisting}

在 \texttt{main.tex} 中，参考文献通过以下方式引入：

\begin{lstlisting}
\usepackage{kaobiblioCJKsc}          % 加载 biblatex + 中文适配
\addbibresource{main.bib}            % 添加 .bib 文献库
% ... 文档主体 ...
\printbibliography[heading=bibintoc] % 打印参考文献列表
\end{lstlisting}

\index{引用!.bib文件@\texttt{.bib} 文件}
参考文献数据存放在 \texttt{.bib} 文件中，每条记录形如：

\begin{lstlisting}
@article{Author2024,
  author  = {张三 and 李四},
  title   = {关于××的研究},
  journal = {××学报},
  year    = {2024},
  volume  = {42},
  pages   = {1--10},
}
\end{lstlisting}

%----------------------------------------------------------------------------------------

\section{索引条目}

\index{索引}
\index{引用!索引}
本模板使用 \Package{imakeidx} 生成索引。在正文中插入
\lstinline|\index{关键词}| 即可将相应条目加入索引。
编译索引需要额外运行 \texttt{makeindex}（\Package{imakeidx} 可在编译时
自动完成这一步骤）。

至此，你应该已经掌握了本模板提供的引用体系：文献引用（\Package{biblatex}）、
交叉引用（\Package{kaorefsCJKsc}）和智能引用（\Package{cleveref}）三管齐下，
能够满足绝大多数中文学术写作的需求。
